\section{Introduction}\label{n:introduction}

Let $(M^n,g)$ be a connected complete smooth Riemannian manifold without
boundary. For $d\geq0$, write $\mathcal H_d(M,g)$ for the real vector space
of smooth harmonic functions such that
\[
 |u(x)|\leq C_u(1+d_g(o,x))^d\qquad(x\in M)
\]
with some finite constant $C_u$ depending on $u$. The choice of base point
$o$ does not change the space. Let $h_d(M,g)$ denote its real dimension.
For an integer $k\geq0$, $\mathcal H_k(\mathbb R^3,g_{\mathrm E})$ consists of harmonic
polynomials of degree at most $k$; summing their homogeneous dimensions gives
\begin{equation}\label{n:euclidean-count}
 h_k(\mathbb R^3,g_{\mathrm E})=\sum_{l=0}^k(2l+1)=(k+1)^2.
\end{equation}
The integer-degree comparison problem asks whether nonnegative Ricci
curvature forces
\[
 h_k(M^n,g)\leq h_k(\mathbb R^n,g_{\mathrm E})
 \qquad\text{for every integer }k\geq1.
\]
Our answer in dimension three is negative.
We use the normalization
\[
 \operatorname{AVR}(g)=\lim_{R\to\infty}\frac{\operatorname{vol}_g B_g(o,R)}{\omega_nR^n},
\]
where $\omega_n$ is the Euclidean unit-ball volume. A pointed tangent cone
at infinity is a pointed Gromov--Hausdorff limit of rescalings
$(M,R_j^{-2}g,o)$ with $R_j\to\infty$.

\begin{theorem}\label{n:main}
For every \(v\in(4/9,1)\) and \(c\in(1,9v/4)\), there is an integer
\(k_0(v,c)\) such that for every integer \(k\geq k_0(v,c)\) there is a
complete smooth metric \(g_k\) on \(\mathbb R^3\), Euclidean near the
origin, with \(\operatorname{Ric}_{g_k}\geq0\),
\(\operatorname{AVR}(g_k)=v\), and
\[
 h_k(\mathbb R^3,g_k)\geq c(k+1)^2>(k+1)^2
 =h_k(\mathbb R^3,g_{\mathrm E}).
\]
The Ricci curvature is positive outside a compact set. The metric has
uncountably many pairwise nonisometric pointed tangent cones at infinity.
There are points tending to infinity and planes containing the radial
direction at those points whose sectional curvature is negative.
The metric is chosen after \(k\) and may depend on it.
\end{theorem}

The excess can occur even when the identity map to Euclidean space has
arbitrarily small bi-Lipschitz distortion.

\begin{corollary}\label{n:near-euclidean}
For every \(\epsilon>0\) and \(1<c<9/4\), all sufficiently large integers
\(k\) admit metrics with the conclusions of Theorem~\ref{n:main} and
\[
 (1+\epsilon)^{-2}g_{\mathrm E}\leq g_k
       \leq(1+\epsilon)^2g_{\mathrm E}
 \qquad\text{on }\mathbb R^3.
\]
Their asymptotic volume ratio can be prescribed arbitrarily close to \(1\).
In particular, the identity map is globally \((1+\epsilon)\)-bi-Lipschitz.
\end{corollary}

The range \(c<9v/4\) comes from the angular average in this construction;
no optimality or endpoint assertion is made. Both statements concern a
family of metrics indexed by the degree, rather than an asymptotic
harmonic-dimension bound on one fixed manifold.

Yau posed the comparison together with the finite-dimensionality problem
for polynomial-growth harmonic functions~\cite{Yau,LiTam}. Li and Tam
proved the sharp bound at linear growth~\cite{LiTam}. Colding and Minicozzi
then proved finite dimensionality at every fixed degree under nonnegative
Ricci curvature and obtained dimension bounds with the optimal order in the
degree~\cite{ColdingMinicozzi97,ColdingMinicozzi98}. An exact count at a
finite integer remains a different question. Donnelly's examples in
dimensions at least five violate the Euclidean comparison at a suitable
real degree between one and two~\cite{Donnelly,CaiLai}. This subquadratic
excess does not by itself establish an integer-degree violation.

A separate construction on odd-dimensional Berger sphere links gives an
integer-degree counterexample in some even ambient dimension at least
eight~\cite[Theorem~1.1 and Section~7]{OpenAIBerger}. That proof works in fixed invariant
spaces of round harmonics and tunes exact transfers using simultaneous real
linear control. The present proof answers the three-dimensional question by
a different route. Its angular eigenspaces move, so it must control leakage
into infinitely many higher modes as well as the exchange of the selected
low modes. No theorem from the Berger-link construction is used in the proof
of Theorem~\ref{n:main}.

Additional curvature or asymptotic hypotheses still yield comparison
results. Cai and Lai prove the Euclidean bound in the locally conformally
flat nonnegative-Ricci setting~\cite{CaiLai}. Lin, Wang and Xu prove the
three-dimensional integer bound under nonnegative sectional curvature and
positive asymptotic volume ratio~\cite[Theorem~1.12]{LinWangXu}. For complete
manifolds with nonnegative Ricci curvature, maximal volume growth, and a
unique tangent cone, Huang's Theorem~1.6 bounds harmonic dimensions above by
the spectral counting function of the cone link~\cite{Huang}. Xu's
three-circles theorem permits nonunique cones under conic-measure hypotheses
when the tested exponent avoids the union of their cone-degree
spectra~\cite[Theorem~3.2]{Xu}; it is not an exact dimension count.
The metrics constructed here have positive asymptotic volume ratio and
distinct cone limits. They also have negative radial sectional planes at
arbitrarily large distances, so they do not meet the sectional-curvature
hypothesis.

\subsection*{Averaging rates by crossing eigenlines}

For a fixed metric $h$ on $S^2$, an eigenvalue $\lambda$ of its
nonnegative Laplacian gives the cone frequency
\[
 d(\lambda)=\frac{-1+\sqrt{1+4\lambda}}2.
\]
An eigenfunction with that frequency grows as $r^{d(\lambda)}$ on the cone.
Our construction does not find one frozen link with too many frequencies
below $k$. In fact, each limiting link has Gauss curvature
greater than one and satisfies the fixed-sphere eigenvalue comparison of
Lin, Wang and Xu~\cite[Theorem~1.8]{LinWangXu}. Instead, a harmonic function
will encounter different frequencies at different radii.

We construct a periodic path $H(s)$ of sphere metrics near the round metric,
all with the same pointwise area form. Its low spectrum is simple except at
finitely many isolated linearly split double crossings. There are two ways to
follow an eigenvalue near such a crossing: reorder by size on each side, or
continue its smooth eigenline through the equality. The latter continuation
exchanges two adjacent ordered positions. A sequence of these exchanges
cycles a long band among the first \(p=(M+1)^2\) positions, where
\(M=\lfloor\vartheta k\rfloor\). For the prescribed \(v,c\), we take
\(a=\sqrt v\) and choose \(\sqrt c<\vartheta<3a/2\). The value of \(p\)
is the round-sphere count through degree \(M\), and
\(p/(k+1)^2\to\vartheta^2>c\).

Over a full cycle of labels, each selected band label visits every position
in the band at the same phases. Its mean frequency is therefore an average
of ordered frequencies, even though its instantaneous frequency may exceed
$k$. The band begins just above round degree \(L=\lfloor\sqrt k\rfloor\). A direct
round-sphere count gives mean degree asymptotic to \(2\vartheta k/3\).
The radial scaling multiplies this by \(a^{-1}\); since
\(2\vartheta/(3a)<1\), a sufficiently small angular distortion leaves
every selected mean strictly below \(k\). The uncycled lower positions
have frequencies \(O(\sqrt k)\). The lower cutoff tends to infinity so
that the required crossings between consecutive degree blocks are
available, while its ratio to \(k\) tends to zero.

The angular construction must prescribe the crossings while excluding all
unwanted coincidences through position $p+1$. Conformal first variations on
a two-dimensional eigenspace provide both trace-free matrix directions.
Exact-multiplicity charts and a finite pool of variations turn this local
fact into an avoidance argument for every possible multiplicity stratum.
A separate step keeps a prescribed double while splitting all other
multiplicities: if every double-preserving variation failed to split
another eigenspace, the two eigenspaces would have common nodal domains,
contradicting their distinct eigenvalues. For the round sphere, the
symmetric-square multiplication property behind
these variations appears in work of Colin de Verdi\`ere and is proved
explicitly by Greilhuber and Kepplinger~\cite{ColinDeVerdiere,GreilhuberKepplinger}.
We also need exchanges between consecutive round degree blocks; an even
perturbation separates parities in opposite orders on the two sides of a
chosen perturbation. The angular section proves that these ingredients
produce the required cyclic program and a uniform gap above position $p$.

\subsection*{Exact harmonic continuation along the program}

Put $t=\log r$ and let the phase move at speed $\eta(t)=t^{-3/4}$. Small
independent controls are placed near the crossings. The polar metric has
the form
\[
 g=dr^2+f(r)^2H_*(\log r),\qquad f(r)/r\longrightarrow a=\sqrt v.
\]
The radial construction is uniform over all control sequences. Its concave
tail supplies radial Ricci curvature of order $\eta^2/r^2$. The mixed Ricci
block is of order $\eta/r^2$, but its loss in the Schur complement is
quadratic because the tangential block of $r^2\operatorname{Ric}$ has a
fixed positive lower bound.
The tail coefficient is chosen to absorb that loss before the radial start
is moved outward. Moser's volume-form method supplies smooth coordinates
with fixed area form~\cite{Moser}; slow-link constructions with a common volume form
also occur in Colding--Naber~\cite{ColdingNaber}. The curvature argument
here checks the particular controlled path and its uniform parameter bounds.

Following a reference eigenline is not yet following an entire harmonic
function. The crossing controls also perturb the instantaneous eigenspaces.
To impose regularity at the center, we use the exact inward map
that sends data on one sphere to the trace of its harmonic extension across
the whole enclosed ball on the next inner sphere. These maps preserve
constants and the spherical mean, are positive, and contract $L^2$. Freezing
the geometry on a long terminal part of the ball gives two estimates with
different purposes: an operator estimate separates the first $p$ modes from
the tail, while a sharper estimate on the finite moving frame detects the
small signed effect of each crossing control.

The outer spectral gap produces a graph over the first $p$ angular modes
that is invariant under the exact inward maps. Its high-mode component
depends on future controls. The induced outward maps on the low coordinates
are finite matrices, and a crossing control changes the relevant
off-diagonal entry with a fixed nonzero first-order sign. To make an invariant
line follow a continued eigenbranch, we solve a scalar recurrence forward
on its contracting side and backward on its other side. Their discrepancy at
the crossing is a positive-weight sum. The two endpoint choices of the
control give opposite signs. Finite-dimensional fixed points followed by a
compact diagonal limit select all crossing controls simultaneously.

For the resulting single metric $g_k$, the matched lines give compatible
boundary traces on nested balls. Whole-ball harmonic extensions then define
$p$ independent smooth harmonic functions on all of $\mathbb R^3$.
Their logarithmic growth is the integral of the selected cone frequencies,
up to errors smaller than the logarithmic radius. Periodic averaging gives
the strict sub-$k$ exponent. The contraction of inward maps controls every
intermediate sphere, and uniform interior elliptic estimates convert the
spherical $L^2$ bounds to the pointwise condition defining $\mathcal H_k$.
This last passage is needed for the integer comparison: the strict average
bound at the sampled radii alone would not suffice.

All functions and Hilbert spaces below are real. Eigenvalues are numbered
from $1$, with the constant eigenfunction at position $1$ and eigenvalue
$0$. The integer $k$, the finite angular program, and $p$ are fixed before
the tail coefficient and the radial start are chosen. Every error estimate
in the transfer and matching argument is for this fixed finite program;
no uniformity as $k\to\infty$ is asserted.


Section~\ref{n:sec-angular} constructs the isolated crossings, and
Section~\ref{n:sec-angular-program} assembles their cycle and proves the
mean-frequency margin. Section~\ref{n:sec-geometry-input} realizes the
controlled links by complete Ricci-nonnegative metrics.
Section~\ref{n:sec-transfer} reduces whole-ball continuation to exact
finite matrices. Section~\ref{n:sec-matching-growth} chooses the crossing
controls, and Section~\ref{n:sec-growth} constructs the entire functions
and proves their growth at every point.
